\documentclass[12pt,reqno]{amsart}
\usepackage{amsmath,amssymb,amsthm}
\usepackage{mathptmx}
\usepackage[scaled=0.90]{helvet}
\usepackage{courier}
\usepackage[T1]{fontenc}
\usepackage{enumitem}
\usepackage[hidelinks]{hyperref}
\pagestyle{plain}

\setlength{\oddsidemargin}{0cm}
\setlength{\evensidemargin}{0cm}
\setlength{\topmargin}{0in}
\setlength{\headheight}{0pt}
\setlength{\headsep}{0pt}
\setlength{\footskip}{.3in}
\setlength{\textheight}{9.2in}
\setlength{\textwidth}{6.5in}
\setlength{\parindent}{0pt}
\setlist[enumerate,1]{leftmargin=*,itemsep=1.1em}
\setlist[enumerate,2]{label=(\alph*),leftmargin=1.7em,itemsep=.45em,topsep=.4em}

\newcommand{\C}{\mathbb{C}}
\newcommand{\R}{\mathbb{R}}

\title{Problem Set 3: Solutions \\ MATH 411: Honors Complex Variables, Fall 2026}
\date{last compiled: \today}

% Coverage: material excluded from Midterm 1 through Monday, Oct. 26,
% 2026 (Lectures 13--20), according to the current course schedule.
% Suggested completion date: Monday, October 26, 2026.
% Sources (paths relative to the course workspace):
% 1: Old Homework/Homework/Fall2021-MATH411-Homework4.tex, 1(a)--(b).
% 2: same source, 2; adapted to the fixed-endpoint definition in Lecture 17.
% 3: same source, 3; f is required to be nonconstant to avoid a trivial
%    counterexample whose image is not open.
% 4: same source, 4--5; argument intervals replace ambiguous branch names.
% 5: same source, 6; sharpened to specified domains and a global obstruction.
% 6: Old Homework/Homework/Homework5_Fall2021.tex, 1.
% Laurent expansions moved to Problem Set 4.
% Residue and winding-number questions from later old sets are omitted.


\usepackage[most]{tcolorbox}
\definecolor{solblue}{RGB}{232,243,252}
\definecolor{solframe}{RGB}{120,170,215}
\newenvironment{solution}
  {\par\smallskip\begin{tcolorbox}[breakable,enhanced,colback=solblue,colframe=solframe,
   boxrule=0.5pt,arc=2pt,left=6pt,right=6pt,top=4pt,bottom=4pt]\noindent\textbf{Solution.}}
  {\hfill$\square$\end{tcolorbox}\par\medskip}
\usepackage{needspace}
\widowpenalty=10000
\clubpenalty=10000
\raggedbottom
% Solutions 2--5 follow Fall2021-MATH411-Homework4_solutions.tex (2--6).
% Solution 6 follows Homework5_Fall2021solutions.pdf (1),
% also present in hwk3.tex.
% Solution 1 supplies the contour arguments requested in the old set,
% which were not written out in the corresponding old solutions.

\begin{document}
\maketitle

Throughout, homotopies of paths \emph{fix their endpoints}, and a simply connected open set is required to be nonempty and connected.

\begin{enumerate}
\item \textbf{Real integrals by contour integration.}
Use Cauchy's theorem to evaluate the following improper integrals.
\begin{enumerate}
\item
\[
 \int_0^\infty \cos(t^2)\,dt.
\]
\emph{Hint:} Integrate $e^{-z^2}$ around the boundary of the sector
$0\leq\arg z\leq\pi/4$, $|z|\leq R$.
You may use
\[
 \int_{-\infty}^{\infty}e^{-x^2}\,dx=\sqrt{\pi},\qquad
 \sin u\geq\frac{2u}{\pi}\quad(0\leq u\leq\pi/2).
\]
Show that the integral over the circular arc tends to zero; a bound
using only the maximum of the integrand on the arc will not suffice.
\item For $a,b>0$, evaluate
\[
 \int_0^\infty e^{-ax}\cos(bx)\,dx.
\]
\emph{Hint:} Put $A=\sqrt{a^2+b^2}$ and choose
$0<\theta<\pi/2$ with $\cos\theta=a/A$ and $\sin\theta=b/A$.
Integrate $e^{-Az}$ around the sector of radius $R$ and angle $\theta$.
\end{enumerate}

\begin{solution}
\noindent (a) Let $C_R$ be the positively oriented boundary of the sector
$0\leq\arg z\leq\pi/4$, $|z|\leq R$, and let $A_R$ be its circular
arc, traversed from $R$ to $Re^{i\pi/4}$. Since $e^{-z^2}$ is entire,
Cauchy's theorem gives
\[
0=\int_0^R e^{-x^2}\,dx+\int_{A_R}e^{-z^2}\,dz
  -e^{i\pi/4}\int_0^R e^{-it^2}\,dt.
\]
We first show that the arc integral approaches zero. Parametrizing by
$z=Re^{i\phi}$, $0\leq\phi\leq\pi/4$, we obtain
\[
\left|\int_{A_R}e^{-z^2}\,dz\right|
 \leq R\int_0^{\pi/4}e^{-R^2\cos(2\phi)}\,d\phi.
\]
Put $u=\pi/4-\phi$. Then $\cos(2\phi)=\sin(2u)\geq4u/\pi$, by the
inequality in the question. Hence
\begin{align*}
\left|\int_{A_R}e^{-z^2}\,dz\right|
 &\leq R\int_0^{\pi/4}e^{-4R^2u/\pi}\,du\\
 &=\frac{\pi}{4R}(1-e^{-R^2})\leq\frac{\pi}{4R}\longrightarrow0.
\end{align*}
Solving the contour identity for the integral on the sloping ray gives
\[
\int_0^R e^{-it^2}\,dt
 =e^{-i\pi/4}\left(\int_0^R e^{-x^2}\,dx
                  +\int_{A_R}e^{-z^2}\,dz\right).
\]
The right hand side has a limit as $R\to\infty$, so this also proves
convergence of the improper integral on the left. Since
$\int_0^\infty e^{-x^2}\,dx=\sqrt{\pi}/2$, taking real parts yields
\[
\boxed{\displaystyle\int_0^\infty\cos(t^2)\,dt
 =\frac{\sqrt{\pi}}{2}\cos(\pi/4)
 =\frac{\sqrt{\pi}}{2\sqrt2}=\sqrt{\frac{\pi}{8}}.}
\]

\medskip
\noindent (b) Put $A=\sqrt{a^2+b^2}$ and choose $\theta$ as in the
hint, so that $Ae^{i\theta}=a+ib$. Integrating $e^{-Az}$ around the
positively oriented sector of radius $R$ and angle $\theta$ gives
\[
0=\int_0^R e^{-Ax}\,dx+\int_{B_R}e^{-Az}\,dz
       -e^{i\theta}\int_0^R e^{-(a+ib)x}\,dx,
\]
where $B_R$ is the circular arc. On this arc, $z=Re^{i\phi}$ with
$0\leq\phi\leq\theta$, so that
\[
|e^{-Az}|=e^{-AR\cos\phi}\leq e^{-AR\cos\theta}=e^{-aR}.
\]
Consequently
\[
\left|\int_{B_R}e^{-Az}\,dz\right|
 \leq R\theta e^{-aR}\longrightarrow0.
\]
Also $|e^{-(a+ib)x}|=e^{-ax}$ is integrable on $[0,\infty)$, so
all the ray integrals converge absolutely. Letting $R\to\infty$,
we obtain
\[
\int_0^\infty e^{-(a+ib)x}\,dx
 =\frac{e^{-i\theta}}{A}=\frac{a-ib}{a^2+b^2}.
\]
Taking real parts now gives
\[
\boxed{\displaystyle\int_0^\infty e^{-ax}\cos(bx)\,dx
       =\frac{a}{a^2+b^2}.}
\]
\end{solution}

\item \textbf{Star-shaped open sets.}
A nonempty open set $U\subseteq\C$ is \emph{star-shaped with respect
to $a\in U$} if the segment $[a,z]$ lies in $U$ for every $z\in U$.
\begin{enumerate}
\item Prove that $U$ is path connected.
\item For a closed continuous curve $\gamma:[0,1]\to U$ based at $a$,
construct a homotopy with fixed endpoints from $\gamma$ to the constant
curve at $a$. Deduce that $U$ is simply connected.
\end{enumerate}
For the last deduction, you may use the elementary fact that in a path
connected set it suffices to contract loops at one fixed base point.
A connecting path and its reverse transfer the contraction to any
other base point.

\begin{solution}
\noindent (a) For any $z,w\in U$, follow the segment from $z$ to $a$
and then the segment from $a$ to $w$. Both segments lie in $U$ by the
star-shaped hypothesis. More explicitly, the continuous path
\[
p(t)=\begin{cases}
 (1-2t)z+2ta,&0\leq t\leq1/2,\\
 (2-2t)a+(2t-1)w,&1/2\leq t\leq1
\end{cases}
\]
joins $z$ to $w$ in $U$. Hence $U$ is path connected.

\medskip
\noindent (b) Let $\gamma:[0,1]\to U$ be a closed curve based at $a$.
Define $F:[0,1]\times[0,1]\to U$ by
\[
F(t,u)=(1-u)\gamma(t)+ua.
\]
Clearly $F(t,u)\in U$ by the star-shaped hypothesis, and $F$ is
continuous. Moreover,
\[
F(t,0)=\gamma(t),\qquad F(t,1)=a,
\]
and, since $\gamma(0)=\gamma(1)=a$,
\[
F(0,u)=F(1,u)=a\qquad(0\leq u\leq1).
\]
Hence $F$ is the required homotopy with fixed endpoints from $\gamma$
to the constant curve at $a$. By (a) and the base-point fact stated in
the question, $U$ is simply connected.
\end{solution}

\item \textbf{Images of simply connected sets.}
Let $U$ be a simply connected open set and let $f$ be a nonconstant
holomorphic function on $U$. Must $f(U)$ be simply connected?
Investigate this using
\[
 H=\{z\in\C:\operatorname{Im}z>0\},\qquad f(z)=e^{2\pi i z}.
\]
Determine $f(H)$ exactly and justify your conclusion about its simple
connectivity, for example by integrating $1/z$ around a suitable circle.

\begin{solution}
The answer is \emph{not necessarily}. The upper half-plane $H$ is
convex, hence star-shaped and simply connected by Question~2.
The function $f(z)=e^{2\pi iz}$ is nonconstant and holomorphic on $H$.

If $z=x+iy$, then
\[
f(z)=e^{-2\pi y}e^{2\pi ix},
\]
the polar decomposition of $f(z)$. The radius of $f(z)$ is thus
between $0$ and $1$ and the argument can be arbitrary. More precisely,
if $w=re^{i\phi}$ with $0<r<1$, take
\[
x=\frac{\phi}{2\pi},\qquad y=-\frac{\log r}{2\pi}>0.
\]
Then $f(x+iy)=w$. Thus
\[
\boxed{f(H)=\{w\in\C:0<|w|<1\},}
\]
the open unit disc minus the origin.

To see that this image is not simply connected, fix $0<r<1$ and
let $\gamma(t)=re^{it}$ for $0\leq t\leq2\pi$. The function $1/z$
is holomorphic on $f(H)$, but
\[
\int_\gamma\frac{dz}{z}
 =\int_0^{2\pi}\frac{ire^{it}}{re^{it}}\,dt=2\pi i\neq0.
\]
If $f(H)$ were simply connected, Cauchy's theorem would force this
integral to be zero, a contradiction.
\end{solution}

\newpage
\item \textbf{Two branches of the logarithm.}
Let $L_+$ be the logarithm on $\C\setminus[0,\infty)$ with argument
in $(0,2\pi)$, and let $L_-$ be the logarithm on
$\C\setminus(-\infty,0]$ with argument in $(-\pi,\pi)$.
For \emph{each} branch $L$, compute
\[
 L(-1+i),\qquad \exp\bigl(iL(-i)\bigr),\qquad L(-1-i).
\]
Here $(-i)^i$ is defined relative to the chosen branch as
$\exp(iL(-i))$. Explain how $L_+$ and $L_-$ differ on the upper and
lower half-planes.

\begin{solution}
For either branch we use $L(z)=\log|z|+i\arg z$ with the specified
interval for the argument.

\medskip
\noindent For $L_+$ the arguments of $-1+i$, $-i$, and $-1-i$ are,
respectively, $3\pi/4$, $3\pi/2$, and $5\pi/4$. Thus
\begin{align*}
 L_+(-1+i)&=\tfrac12\log2+\tfrac{3\pi i}{4},\\
 \exp(iL_+(-i))&=\exp\bigl(i\cdot\tfrac{3\pi i}{2}\bigr)
                  =e^{-3\pi/2},\\
 L_+(-1-i)&=\tfrac12\log2+\tfrac{5\pi i}{4}.
\end{align*}

\noindent For $L_-$ the corresponding arguments are $3\pi/4$,
$-\pi/2$, and $-3\pi/4$. Hence
\begin{align*}
 L_-(-1+i)&=\tfrac12\log2+\tfrac{3\pi i}{4},\\
 \exp(iL_-(-i))&=\exp\bigl(i\cdot(-\tfrac{\pi i}{2})\bigr)
                  =e^{\pi/2},\\
 L_-(-1-i)&=\tfrac12\log2-\tfrac{3\pi i}{4}.
\end{align*}
The branches have the same argument on the upper half-plane.
On the lower half-plane the argument for $L_+$ is $2\pi$ greater
than the one for $L_-$. Therefore
\[
L_+(z)-L_-(z)=
\begin{cases}
0,&\operatorname{Im}z>0,\\
2\pi i,&\operatorname{Im}z<0.
\end{cases}
\]
\end{solution}

\item \textbf{Holomorphic roots and the choice of domain.}
Fix an integer $n\geq2$.
\begin{enumerate}
\item Let $U\subseteq\C\setminus\{0\}$ be a simply connected open
set. Use a holomorphic logarithm on $U$ to construct a holomorphic
function $g$ satisfying $g(z)^n=z$ on $U$.
\item For $\theta\in\R$, give such a branch on the slit plane
\[
 U_\theta=\C\setminus\{re^{i\theta}:r\geq0\},
\]
specifying the interval used for the argument.
\item Show that there is no continuous function $g$ on
$\C\setminus\{0\}$ satisfying $g(z)^n=z$ everywhere.
\emph{Hint:} On the unit circle, consider
$g(e^{it})e^{-it/n}$ for $0\leq t\leq2\pi$.
\end{enumerate}

\needspace{22\baselineskip}
\begin{solution}
\noindent (a) On any simply connected open set not containing the
origin, a holomorphic logarithm $L$ is well-defined, with $e^{L(z)}=z$.
Define
\[
\boxed{g(z)=\exp\left(\frac{L(z)}{n}\right).}
\]
This is holomorphic on $U$, and $g(z)^n=e^{L(z)}=z$.

\medskip
\noindent (b) On the slit plane $U_\theta$, choose the continuous
argument $\arg_\theta z\in(\theta,\theta+2\pi)$. Then
\[
L_\theta(z)=\log|z|+i\arg_\theta z
\]
is a holomorphic logarithm, and a branch of the $n$th root is
\[
\boxed{g_\theta(z)=|z|^{1/n}
                 \exp\left(\frac{i\arg_\theta z}{n}\right).}
\]
For example, deleting the negative real axis allows the argument
interval $(-\pi,\pi)$ (take $\theta=-\pi$).

\medskip
\noindent (c) Suppose there were a continuous function $g$ on
$\C\setminus\{0\}$ satisfying $g(z)^n=z$. On $[0,2\pi]$ put
\[
h(t)=g(e^{it})e^{-it/n}.
\]
Then $h$ is continuous and
\[
h(t)^n=g(e^{it})^n e^{-it}=1.
\]
Thus $h$ takes its values in the finite set of $n$th roots of unity.
Since $[0,2\pi]$ is connected, $h$ must be constant, say $h(t)=c$
with $c^n=1$. It follows that
\[
g(e^{it})=ce^{it/n}.
\]
Evaluating at $t=0$ and $t=2\pi$ gives
\[
g(1)=c\qquad\text{and}\qquad g(1)=ce^{2\pi i/n}.
\]
Since $c\neq0$ and $e^{2\pi i/n}\neq1$ for $n\geq2$, this is a
contradiction. Hence no such continuous function exists on the entire
punctured plane.
\end{solution}

\needspace{28\baselineskip}
\item \textbf{Polynomial growth and Cauchy estimates.}
Let $f$ be entire and put
\[
 M_f(R)=\max_{|z|=R}|f(z)|.
\]
Suppose there are a constant $C>0$, an integer $k\geq0$, and radii
$R_j\to\infty$ such that $M_f(R_j)\leq C R_j^k$ for every $j$.
Use Cauchy's estimates for the Taylor coefficients at $0$ to prove
that $f$ is a polynomial of degree at most $k$ (including the zero
polynomial). What does the case $k=0$ say?

\begin{solution}
Write the Taylor expansion of the entire function as
\[
f(z)=\sum_{m=0}^\infty a_mz^m.
\]
Let $N>k$. Then using the Cauchy estimates on the coefficients in
this power series expansion, we find for each positive radius $R_j$
\[
|a_N|\leq\frac{M_f(R_j)}{R_j^N}
       \leq C\frac{R_j^k}{R_j^N}=\frac{C}{R_j^{N-k}}.
\]
As $R_j\to\infty$ the right hand side tends to zero, so $a_N=0$.
This holds for every $N>k$, i.e.\ in the power series expansion of
$f$, only finitely many coefficients can be nonzero. Hence
\[
f(z)=a_0+a_1z+\cdots+a_kz^k,
\]
a polynomial of degree at most $k$, including the possibility that
$f$ is the zero polynomial.

For $k=0$ all coefficients except $a_0$ vanish, so $f$ is constant.
This is the conclusion of Liouville's theorem, here obtained from a
uniform bound on circles of arbitrarily large radius.
\end{solution}

\end{enumerate}
\end{document}
